\documentclass[11pt,a4paper]{article}

% --- Packages ---
\usepackage[utf8]{inputenc}
\usepackage[T1]{fontenc}
\usepackage[margin=2.5cm]{geometry}
\usepackage{amsmath,amssymb}
\usepackage{booktabs}
\usepackage{enumitem}
\usepackage{hyperref}
\usepackage{xcolor}
\usepackage{fancyhdr}
\usepackage{parskip}
\usepackage{graphicx}
\usepackage{float}

% --- Colors ---
\definecolor{darkblue}{HTML}{1a1a2e}
\definecolor{accentblue}{HTML}{0f3460}
\definecolor{lightgray}{HTML}{f5f5f5}

% --- Hyperlinks ---
\hypersetup{
    colorlinks=true,
    linkcolor=accentblue,
    urlcolor=accentblue,
    citecolor=accentblue,
}

% --- Section colors ---
\makeatletter
\renewcommand\section{\@startsection{section}{1}{\z@}{-3.5ex \@plus -1ex \@minus -.2ex}{2.3ex \@plus.2ex}{\normalfont\large\bfseries\color{darkblue}}}
\renewcommand\subsection{\@startsection{subsection}{2}{\z@}{-3.25ex\@plus -1ex \@minus -.2ex}{1.5ex \@plus .2ex}{\normalfont\normalsize\bfseries\color{accentblue}}}
\makeatother

% --- Header/Footer ---
\pagestyle{fancy}
\fancyhf{}
\fancyhead[L]{\small\textcolor{gray}{Commodity Return Forecasting}}
\fancyhead[R]{\small\textcolor{gray}{Nassim Taier}}
\fancyfoot[C]{\small\thepage}
\renewcommand{\headrulewidth}{0.4pt}

\begin{document}

% ============================================================
% TITLE PAGE
% ============================================================
\begin{titlepage}
\centering
\vspace*{4cm}

{\Huge\bfseries\color{darkblue} Commodity Return Forecasting}\\[0.8cm]
{\Large Predicting 1-Week and 1-Month Returns\\Using Machine Learning}\\[2cm]

\textcolor{gray}{\rule{8cm}{0.5pt}}\\[1.5cm]

{\large Nassim Taier}\\[0.4cm]
{\normalsize M2 Probability and Finance (El Karoui)\\
Sorbonne Universit\'e \& \'Ecole Polytechnique}\\[1.5cm]

{\normalsize March 2026}

\vfill
\end{titlepage}

% ============================================================
% TABLE OF CONTENTS
% ============================================================
\tableofcontents
\newpage

% ============================================================
% 1. INTRODUCTION
% ============================================================
\section{Introduction}

This report presents a systematic approach to forecasting short-term commodity returns using machine learning. The objective is to predict 1-week and 1-month returns from 60 anonymised features, with a strict methodology designed to prevent data leakage and produce realistic, out-of-sample performance estimates.

The key methodological principles are:
\begin{enumerate}[nosep]
    \item All preprocessing (imputation, scaling) is fitted on training data only.
    \item Cross-validation uses temporal splits exclusively --- no random K-fold.
    \item Target lags are shifted by the forecast horizon to avoid overlap.
    \item An embargo gap separates train and test sets in walk-forward folds.
    \item All reported metrics are computed on non-overlapping observations.
\end{enumerate}

The analysis progresses from simple linear baselines (ElasticNet) to gradient-boosted trees (XGBoost) with engineered features, demonstrating that careful feature engineering improves signal quality beyond what the model can extract from raw inputs alone.

% ============================================================
% 2. DATA
% ============================================================
\section{Data Description}

The dataset consists of daily observations with 60 anonymised features and two target variables: 1-week forward returns (\texttt{output1\_1w}) and 1-month forward returns (\texttt{output1\_1m}). Features are provided without labels, preventing economic interpretation but allowing a purely statistical approach to signal extraction.

\begin{table}[H]
\centering
\caption{Dataset overview}
\begin{tabular}{@{}ll@{}}
\toprule
\textbf{Property} & \textbf{Value} \\
\midrule
Features & 60 (anonymised) \\
Targets & 1-week return, 1-month return \\
Missing values & $\sim$15\% overall (varies by feature) \\
Horizons & 5 trading days (1w), 21 trading days (1m) \\
\bottomrule
\end{tabular}
\end{table}

% ============================================================
% 3. EDA
% ============================================================
\section{Exploratory Data Analysis}

\subsection{Target Distributions}

Both target distributions are approximately symmetric around zero with heavy tails (excess kurtosis), consistent with financial return distributions. The 1-month target exhibits higher variance, as expected for a longer horizon. Augmented Dickey-Fuller tests confirm stationarity of both series ($p < 0.01$).

\subsection{Feature--Target Correlations}

Individual feature--target correlations are generally weak ($|r| < 0.15$), which is expected for financial data: the signal is distributed across many features rather than concentrated in a few. This motivates the use of ensemble tree models that can combine many weak predictors.

\subsection{Predictiveness Stability}

Rolling correlation analysis (252-day window) reveals that feature predictiveness is not stable over time. Features that are strong predictors in one period may lose predictive power in another. This non-stationarity motivates walk-forward evaluation rather than a single train/test split, and suggests that adaptive or regime-conditional models could further improve performance.

\subsection{Feature Redundancy}

The feature correlation matrix reveals clusters of highly correlated features ($|r| > 0.9$). This redundancy is partially addressed by XGBoost's built-in feature selection and by L1 regularisation in ElasticNet. The engineered features (z-scores, ranks) provide alternative representations that reduce sensitivity to multicollinearity.

% ============================================================
% 4. FEATURE ENGINEERING
% ============================================================
\section{Feature Engineering}

A total of 140+ features were engineered from the original 60, using only backward-looking transformations to prevent any form of look-ahead bias:

\begin{itemize}[nosep]
    \item \textbf{Rolling z-scores} (21d, 63d): Normalise each feature by its recent rolling mean and standard deviation, capturing whether current values are unusual relative to the recent regime.
    \begin{equation}
        z_{t}^{(w)} = \frac{x_t - \bar{x}_{t-w:t}}{\sigma_{t-w:t} + \epsilon}
    \end{equation}

    \item \textbf{Cross-sectional percentile ranks}: Rank each feature value relative to all other features on the same date, capturing relative positioning across the feature space. This is robust to outliers and non-stationarity in levels.

    \item \textbf{Rolling volatility}: Standard deviation over 21 days for the top 10 most correlated features, capturing changes in feature stability that may signal regime transitions.

    \item \textbf{Feature lags}: Lagged values of top features, shifted by the forecast horizon $h$ to avoid overlap:
    \begin{equation}
        x_{t}^{\text{lag}_k} = x_{t - (h + k)}, \quad k \in \{0, 1, 2, 5, 10\}
    \end{equation}

    \item \textbf{Momentum features}: Change in feature value over the forecast horizon, measuring the recent trend in each predictor:
    \begin{equation}
        \Delta x_t^{(h)} = x_t - x_{t-h}
    \end{equation}
\end{itemize}

These transformations increase the feature space from 60 to approximately 200 features. XGBoost's internal feature selection via split gain and L2 regularisation handle the increased dimensionality without explicit feature selection.

% ============================================================
% 5. MODELING
% ============================================================
\section{Modeling Approach}

\subsection{Model Progression}

Three model classes are evaluated in a principled progression from simple to complex:

\begin{enumerate}[nosep]
    \item \textbf{ElasticNet}: Linear baseline with combined L1/L2 regularisation. Internal cross-validation uses \texttt{TimeSeriesSplit} with 3 folds. Preprocessing (mean imputation + standard scaling) is fitted on train data only via a scikit-learn \texttt{Pipeline}.

    \item \textbf{XGBoost (raw features)}: Gradient-boosted trees on the original 60 features plus horizon-shifted target lags. XGBoost handles missing values natively, eliminating the need for imputation.

    \item \textbf{XGBoost (engineered features)}: Same architecture on the expanded feature set ($\sim$200 features). This tests whether explicit feature engineering adds predictive power beyond what the tree model discovers through its own splits.
\end{enumerate}

\subsection{Evaluation Framework}

All models are evaluated exclusively through 5-fold walk-forward cross-validation. Each fold uses a growing training window and a fixed-size test window, with an embargo gap of $h$ days (where $h$ is the forecast horizon) between the last training observation and the first test observation:

\begin{equation}
    \text{Train}_k = \{1, \ldots, T_k - h\}, \quad \text{Test}_k = \{T_k + 1, \ldots, T_{k+1}\}
\end{equation}

No single train/test split results are reported, as they are sensitive to the arbitrary choice of split point and can be misleadingly optimistic or pessimistic.

\subsection{Metrics}

Primary evaluation metrics are:
\begin{itemize}[nosep]
    \item \textbf{Pearson IC} (information coefficient): Correlation between predicted and realised returns. Computed on non-overlapping observations only (every $h$-th observation).
    \item \textbf{Hit rate}: Fraction of correct directional predictions (non-overlapping).
    \item \textbf{Annualised Sharpe ratio}: Mean PnL / Std PnL $\times \sqrt{252/h}$, computed on non-overlapping observations after transaction costs.
\end{itemize}

R-squared is deliberately not used as a primary metric, as it is less interpretable for signal-based trading strategies where even small positive correlations can be economically significant.

% ============================================================
% 6. BACKTEST
% ============================================================
\section{Walk-Forward Backtest}

\subsection{Position Sizing}

Positions are sized proportionally to signal strength rather than using a binary sign function:
\begin{equation}
    w_t = \text{clip}\left(\frac{\hat{y}_t}{\hat{\sigma}_{\hat{y}}}, -1, 1\right)
\end{equation}
where $\hat{y}_t$ is the predicted return, $\hat{\sigma}_{\hat{y}}$ is the standard deviation of predictions, and the clip function constrains positions to $[-1, 1]$. This ensures that stronger signals receive larger allocations while capping maximum exposure.

\subsection{Transaction Costs}

A cost of 2 basis points per unit of turnover is applied:
\begin{equation}
    \text{Cost}_t = \frac{2}{10000} \times |w_t - w_{t-1}|
\end{equation}
Net PnL is then $\text{PnL}_t^{\text{net}} = w_t \cdot r_t - \text{Cost}_t$, where $r_t$ is the realised return.

\subsection{Risk Metrics}

Beyond Sharpe ratio, the following risk metrics are reported:
\begin{itemize}[nosep]
    \item \textbf{Maximum drawdown}: Largest peak-to-trough decline in cumulative PnL.
    \item \textbf{Calmar ratio}: Annualised return divided by maximum drawdown.
    \item \textbf{Win/loss ratio}: Average profit on winning trades divided by average loss on losing trades.
    \item \textbf{Average daily turnover}: Mean absolute change in position.
\end{itemize}

% ============================================================
% 7. RESULTS
% ============================================================
\section{Results}

\subsection{Main Results}

\begin{table}[H]
\centering
\caption{Walk-forward results --- XGBoost with engineered features (non-overlapping metrics)}
\begin{tabular}{@{}lcc@{}}
\toprule
\textbf{Metric} & \textbf{1-Week} & \textbf{1-Month} \\
\midrule
Pearson IC & 0.45 & 0.41 \\
Hit Rate & 66--67\% & 63\% \\
Annualised Sharpe & 2.2 & 1.6 \\
Transaction Costs & 2 bps & 2 bps \\
\bottomrule
\end{tabular}
\end{table}

\subsection{Model Comparison}

\begin{table}[H]
\centering
\caption{IC progression across models (approximate, averaged across walk-forward folds)}
\begin{tabular}{@{}lccc@{}}
\toprule
\textbf{Model} & \textbf{IC (1w)} & \textbf{IC (1m)} & \textbf{Improvement} \\
\midrule
ElasticNet & $\sim$0.30 & $\sim$0.25 & Baseline \\
XGBoost (raw features) & $\sim$0.38 & $\sim$0.35 & +25\% \\
XGBoost (engineered) & $\sim$0.45 & $\sim$0.41 & +50\% \\
\bottomrule
\end{tabular}
\end{table}

Feature engineering contributes approximately 15--20\% improvement in IC over raw features, demonstrating that domain-informed transformations add value beyond what the tree model discovers through its own splits. The largest improvements come from rolling z-scores and momentum features.

\subsection{Interpretation}

An information coefficient of 0.45 on the 1-week horizon is strong by the standards of financial return prediction, where ICs above 0.05--0.10 are generally considered exploitable at scale. The hit rate of 66\% confirms directional predictability well above the 50\% random baseline.

The annualised Sharpe ratio of 2.2 (1-week) after transaction costs indicates a high-quality signal. However, these results should be interpreted with caution: they apply to a single commodity with anonymised features, and real-world implementation would face additional challenges including market impact, liquidity constraints, and regime changes.

Performance varies across walk-forward folds, with some folds producing Sharpe ratios below 1.0. This heterogeneity is consistent with regime-dependent predictability and highlights the importance of reporting fold-level statistics rather than aggregate averages alone.

% ============================================================
% 8. FEATURE IMPORTANCE
% ============================================================
\section{Feature Importance Analysis}

SHAP (SHapley Additive exPlanations) values are computed on the last walk-forward fold to understand which features and engineered transformations drive predictions. Unlike split-based importance, SHAP provides both magnitude and direction of each feature's contribution.

Key findings:
\begin{itemize}[nosep]
    \item Rolling z-scores and momentum features appear prominently among the top 20 predictors, confirming the value of feature engineering.
    \item Target lags contribute meaningfully but do not dominate the importance ranking, suggesting the signal is not purely autoregressive.
    \item Several raw features remain important even after engineering, indicating they carry information not fully captured by the transformations.
    \item Interactions between raw features and their z-score variants suggest the model benefits from seeing both level and regime-normalised representations.
\end{itemize}

% ============================================================
% 9. LIMITATIONS
% ============================================================
\section{Limitations and Future Work}

\subsection{Limitations}

\begin{itemize}[nosep]
    \item \textbf{Single asset}: All results apply to one commodity. Cross-asset generalisation is untested; the signal may be commodity-specific.
    \item \textbf{Anonymised features}: The absence of feature labels prevents economic interpretation, making it impossible to assess whether the model captures known risk premia or novel patterns.
    \item \textbf{Regime dependence}: Performance varies meaningfully across walk-forward folds. Residual autocorrelation in some folds indicates unmodelled temporal structure.
    \item \textbf{No hyperparameter optimisation}: XGBoost parameters are set to reasonable defaults but not optimised per fold. Overfitting risk is mitigated by regularisation but not eliminated.
    \item \textbf{No market impact modelling}: The backtest assumes fills at the 2 bps cost level regardless of position size, which may understate real-world execution costs for larger allocations.
\end{itemize}

\subsection{Future Work}

\begin{enumerate}[nosep]
    \item \textbf{Hyperparameter tuning}: Bayesian optimisation (e.g.\ Optuna) on walk-forward folds, with the outer fold providing an unbiased estimate of generalisation.
    \item \textbf{Ensemble methods}: Blend XGBoost and ElasticNet predictions, potentially weighting by inverse fold-level variance.
    \item \textbf{Regime detection}: Use volatility clustering or hidden Markov models to build conditional models that adapt to market regimes.
    \item \textbf{Multi-asset extension}: Test whether the signal and feature engineering pipeline generalise across commodities.
    \item \textbf{Alternative architectures}: LightGBM, CatBoost, and temporal convolutional networks for capturing longer-range dependencies.
\end{enumerate}

% ============================================================
% 10. CONCLUSION
% ============================================================
\section{Conclusion}

This project demonstrates that commodity returns at the 1-week and 1-month horizons contain exploitable predictive signal. A walk-forward validated XGBoost model with engineered features achieves an information coefficient of 0.45 (1-week) and 0.41 (1-month), with hit rates above 63\% and annualised Sharpe ratios of 1.6 to 2.2 after transaction costs.

The central methodological contribution is the anti-leakage discipline applied throughout. Without temporal splits, horizon-shifted lags, embargo gaps, and non-overlapping metrics, the same data can produce dramatically inflated performance numbers. Every design choice in this pipeline was made to ensure that reported results would be reproducible in a live trading environment.

Feature engineering contributes meaningfully: rolling z-scores, cross-sectional ranks, and momentum features improve IC by 15--20\% over raw inputs. The progression from linear to tree-based models shows diminishing but real returns to model complexity. The signal is real but regime-dependent, and any production deployment would require continuous monitoring and periodic recalibration.

\end{document}
